\section{Related Work}
\label{sec:related}

\subsection{Diffusion Models for Image Synthesis}
Denoising Diffusion Probabilistic Models (DDPMs) and Latent Diffusion Models (LDMs)~\cite{rombach2022highresolution} generate high-fidelity images by iteratively reversing a forward Gaussian diffusion process. Recently, cascaded architectures such as Stable Cascade (W{\"u}rstchen)~\cite{pernias2023stablecascade} introduced extreme latent compression via a multi-stage pipeline: Stage C operates on a highly compressed latent space ($42\times$ spatial compression) to capture global semantics, while Stage B and Stage A decode latents back to high-resolution pixel space. This hierarchical separation provides an ideal substrate for disentangling high-level semantic geometry from low-level stylistic textures.

\subsection{Training-Based Style Personalization}
Early personalization techniques adapted pre-trained diffusion models through optimization. DreamBooth~\cite{ruiz2023dreambooth} and Textual Inversion optimize subject embeddings or model weights, while LoRA and Custom Diffusion enable parameter-efficient adaptation. For style-specific tasks, StyleDrop~\cite{sohn2023styledrop} uses iterative feedback fine-tuning, and B-LoRA~\cite{chung2024blora} trains separate LoRA ranks for style and content blocks. However, all training-based methods suffer from high computational overhead (requiring minutes of GPU gradient descent per style) and risk severe overfitting or semantic memorization when provided with only a single reference exemplar.

\subsection{Training-Free Style Transfer and Attention Control}
To circumvent training costs, recent works manipulate the internal attention layers and latent representations of diffusion models at inference time. StyleAligned~\cite{hertz2023stylealigned} enforces shared self-attention across a batch of images to synchronize style, but relies on DDIM inversion which frequently degrades fine spatial details. InstantStyle~\cite{wang2024instantstyle} uses IP-Adapter~\cite{ye2023ipadapter} feature injection into selected layers to isolate style, yet still exhibits noticeable semantic leakage when the style image contains distinct objects. StyleID~\cite{dinh2024styleid} and DeadDiff~\cite{qi2024deaddiff} attempt mathematical decoupling of style and content, but often dilute stylistic intensity to protect content geometry. In contrast, \method achieves strict mathematical decoupling via orthogonal subspace projections directly inside cross-attention layers, eliminating leakage without diminishing style potency.

\subsection{Guidance Mechanisms and Style Representation}
Guidance techniques direct diffusion sampling toward target distributions. Classifier-Free Guidance (CFG) scales conditional score estimates, while external loss guidance adds task-specific gradients into intermediate latents. While RB-Modulation~\cite{rout2024rbmodulation} formulated terminal cost guidance using Content-Style Disentangled (CSD) models, CSD embeddings often suffer from limited domain coverage and require specialized training. Furthermore, unconstrained gradient guidance pushes sampling trajectories off the valid diffusion manifold. \method addresses this by employing self-supervised DINO-ViT representations~\cite{caron2021emerging,oquab2023dinov2}---which inherently capture mid-level texture co-occurrences without semantic entanglement---and enforcing score-orthogonal gradient projection to guarantee manifold stability.
